\scanpage{102}
a) $t_0$ ens de card 3

b) $\widetilde\Gamma$ ens \inserted{$\ne\varnothing$} avec involution $\deleted{\unsure{\sim}}\sigma_{\widetilde\Gamma}$ (sans pts \unsure{fixe|fixes})

c) Une famille $(\overline E_\alpha)_{\alpha\in t_0}$ de quotients de $\widetilde\Gamma$ \deleted{\unsure{par|plus}}

Avec les conditions suivantes

$\alpha$) Les $\widetilde E_\alpha$ stables par $\deleted{\unsure{\sim}}\sigma_{\widetilde\Gamma}$ (\unsure{donc|soit|sur|dont} $\sigma_\alpha$ induit) et les $\sigma_\alpha$ sont sans pts fixes

$\beta$) $\deleted{\unsure{\widetilde\Gamma\in|\widetilde E}}\forall\alpha,\beta\in t_0,\ \alpha\ne\beta$, l'application
\[
\widetilde\Gamma\xrightarrow{\widetilde\varphi_{\alpha,\beta}}\widetilde E_\alpha\times\widetilde E_\beta
\]
est \uline{injective}, et si $\Gamma_{\alpha\beta}$ est son image, \deleted{on a la condition} et si $u\in\unsure{\mathcal P|\mathfrak P}_2(\widetilde E_\alpha)$ \deleted{sont distincts et conjugués par} \inserted{est une orbite sous} $\sigma_\alpha$ \deleted{$(s\ne s')$} et $t\in\widetilde E_\beta$, on a : $\exists!\,s\in u$ tel que $(s,t)\in\widetilde\Gamma_{\alpha,\beta}$

\inserted{\textcircled{$\gamma$} l'axiome \uline{(**) plus} haut sur $(s_\alpha)_{\alpha\in t_0}$ et $(\underline t_\alpha)_{\alpha\in t_0}$}

\deleted{On peut déduire de $\beta$), posant $\Gamma=\widetilde\Gamma/\sigma_{\widetilde\Gamma}$, $E_\alpha=\widetilde E_\alpha/\widetilde\sigma_\alpha$, que}
\[
\deleted{\Gamma\xrightarrow{\varphi_{\alpha\beta}}{}}
\]

Soit \deleted{\unsure{$E?$}} $\Gamma=\widetilde\Gamma/\sigma_{\widetilde\Gamma}$, $E_\alpha=\widetilde E_\alpha/\sigma_\alpha$, donc on a des quotients $E_\alpha$ de $\Gamma$
\[
\Gamma\xrightarrow{\varphi_\alpha} E_\alpha,\qquad\Gamma\xrightarrow{\varphi}\prod_{\alpha\in t_0}E_\alpha
\]
et l'axiome a) implique que \inserted{$\forall\alpha$,} $\widetilde\Gamma$ s'identifie à l'image inverse des revêtmts d'ordre 2 $\overline E_\alpha$ de $E_\alpha$. Donc la famille de données
% archive sheet number: 101
